
\documentclass{article}
\usepackage{colortbl}
\usepackage{makecell}
\usepackage{multirow}
\usepackage{supertabular}

\begin{document}

\newcounter{utterance}

\centering \large Interaction Transcript for game `wordle', experiment `medium\_frequency\_words\_no\_clue\_no\_critic', episode 3 with nemotron{-}3{-}ultra.
\vspace{24pt}

{ \footnotesize  \setcounter{utterance}{1}
\setlength{\tabcolsep}{0pt}
\begin{supertabular}{c@{$\;$}|p{.15\linewidth}@{}p{.15\linewidth}p{.15\linewidth}p{.15\linewidth}p{.15\linewidth}p{.15\linewidth}}
    \# & \multicolumn{2}{c}{Player} && \multicolumn{2}{c}{Game Master} \\
    \hline

    \theutterance \stepcounter{utterance}  
    & & & \multicolumn{4}{p{0.6\linewidth}}{
        \cellcolor[rgb]{0.9,0.9,0.9}{
            \makecell[{{p{\linewidth}}}]{
                \texttt{\tiny{[P1$\langle$GM]}}
                \texttt{Du är en språkmagiker som gillar att gissa ord genom att använda de givna reglerna.} \\
\\ 
\texttt{Välkommen till Wordle! Du har sex försök på dig att gissa målordet, ett giltigt svenskt ord med fem gemena bokstäver. Använd taggarna "förklaring:" och "gissning:" för att ge en kortfattad förklaring till varje gissning.} \\
\\ 
\texttt{Till exempel, om din gissning är "fråga", ska ditt svar vara} \\
\texttt{förklaring: detta är ett vanligt svenskt ord med fem bokstäver, och jag börjar min gissning med detta ord.} \\
\texttt{gissning: fråga} \\
\\ 
\texttt{Efter varje gissning kommer ditt svar att valideras, och du kommer att få återkoppling som visar vilka bokstäver som är korrekta (gröna), vilka bokstäver som är korrekta men på fel position (gula), och vilka bokstäver som är felaktiga (röda). Denna återkoppling kan vara användbar för att avgöra vilka bokstäver du ska ta med eller utesluta i din nästa gissning.} \\
\\ 
\texttt{Till exempel kan återkopplingen för "fråga" vara:} \\
\texttt{gissning\_återkoppling: f<gul> r<gul> å<grön> g<gul> a<röd>} \\
\\ 
\texttt{Förklaringen ska innehålla detaljer om hur gissning\_återkoppling används för att komma fram till en ny gissning.} \\
\\ 
\texttt{Generera ingen annan text.} \\
\\ 
\texttt{Låt oss börja med din första gissning.} \\
            }
        }
    }
     \\ \\

    \theutterance \stepcounter{utterance}  
    & \multicolumn{4}{p{0.6\linewidth}}{
        \cellcolor[rgb]{0.561,0.737,0.561}{
            \makecell[{{p{\linewidth}}}]{
                \texttt{\tiny{[P1$\rangle$GM]}}
                \texttt{förklaring: jag börjar med ett vanligt ord som innehåller tre vokaler och två vanliga konsonanter för att maximera informationen vid första gissningen.} \\
\texttt{gissning: radio} \\
            }
        }
    }
    & & \\ \\

    \theutterance \stepcounter{utterance}  
    & & & \multicolumn{2}{p{0.3\linewidth}}{
        \cellcolor[rgb]{0.9,0.9,0.9}{
            \makecell[{{p{\linewidth}}}]{
                \texttt{\tiny{[GM$|$GM]}}
                \texttt{attempts = 1} \\
\texttt{target\_word = butik} \\
\texttt{guess = radio} \\
\texttt{guess\_feedback = r<red> a<red> d<red> i<green> o<red>} \\
            }
        }
    }
    & & \\ \\

    \theutterance \stepcounter{utterance}  
    & & & \multicolumn{4}{p{0.6\linewidth}}{
        \cellcolor[rgb]{0.9,0.9,0.9}{
            \makecell[{{p{\linewidth}}}]{
                \texttt{\tiny{[P1$\langle$GM]}}
                \texttt{gissning\_återkoppling: r<röd> a<röd> d<röd> i<grön> o<röd>} \\
\\ 
\texttt{Ge ditt svar endast i detta format.} \\
\texttt{förklaring: detaljer} \\
\texttt{gissning: ord} \\
            }
        }
    }
     \\ \\

    \theutterance \stepcounter{utterance}  
    & \multicolumn{4}{p{0.6\linewidth}}{
        \cellcolor[rgb]{0.561,0.737,0.561}{
            \makecell[{{p{\linewidth}}}]{
                \texttt{\tiny{[P1$\rangle$GM]}}
                \texttt{förklaring: bokstaven i är rätt på tredje plats. r, a, d, o finns inte i ordet. jag gissar på ett ord som har i i mitten, slutar på e, och innehåller vanliga konsonanter som inte testats ännu.} \\
\texttt{gissning: siten} \\
            }
        }
    }
    & & \\ \\

    \theutterance \stepcounter{utterance}  
    & & & \multicolumn{2}{p{0.3\linewidth}}{
        \cellcolor[rgb]{0.9,0.9,0.9}{
            \makecell[{{p{\linewidth}}}]{
                \texttt{\tiny{[GM$|$GM]}}
                \texttt{Error: The guessed word is not a valid word for this game.} \\
            }
        }
    }
    & & \\ \\

    \theutterance \stepcounter{utterance}  
    & & & \multicolumn{4}{p{0.6\linewidth}}{
        \cellcolor[rgb]{0.9,0.9,0.9}{
            \makecell[{{p{\linewidth}}}]{
                \texttt{\tiny{[P1$\langle$GM]}}
                \texttt{Din gissning är inte ett giltigt ord för det här spelet. Vänligen försök igen.} \\
\\ 
\texttt{Ge ditt svar endast i detta format.} \\
\texttt{förklaring: detaljer} \\
\texttt{gissning: ord} \\
            }
        }
    }
     \\ \\

    \theutterance \stepcounter{utterance}  
    & \multicolumn{4}{p{0.6\linewidth}}{
        \cellcolor[rgb]{0.561,0.737,0.561}{
            \makecell[{{p{\linewidth}}}]{
                \texttt{\tiny{[P1$\rangle$GM]}}
                \texttt{förklaring: "siten" var inte ett giltigt ord. jag försöker med "sikte" som har i på tredje plats, slutar på e, och använder vanliga bokstäver som inte är uteslutna (s, k, t).} \\
\texttt{gissning: sikte} \\
            }
        }
    }
    & & \\ \\

    \theutterance \stepcounter{utterance}  
    & & & \multicolumn{2}{p{0.3\linewidth}}{
        \cellcolor[rgb]{0.9,0.9,0.9}{
            \makecell[{{p{\linewidth}}}]{
                \texttt{\tiny{[GM$|$GM]}}
                \texttt{attempts = 2} \\
\texttt{target\_word = butik} \\
\texttt{guess = sikte} \\
\texttt{guess\_feedback = s<red> i<yellow> k<yellow> t<yellow> e<red>} \\
            }
        }
    }
    & & \\ \\

    \theutterance \stepcounter{utterance}  
    & & & \multicolumn{4}{p{0.6\linewidth}}{
        \cellcolor[rgb]{0.9,0.9,0.9}{
            \makecell[{{p{\linewidth}}}]{
                \texttt{\tiny{[P1$\langle$GM]}}
                \texttt{gissning\_återkoppling: s<röd> i<gul> k<gul> t<gul> e<röd>} \\
\\ 
\texttt{Ge ditt svar endast i detta format.} \\
\texttt{förklaring: detaljer} \\
\texttt{gissning: ord} \\
            }
        }
    }
     \\ \\

    \theutterance \stepcounter{utterance}  
    & \multicolumn{4}{p{0.6\linewidth}}{
        \cellcolor[rgb]{0.561,0.737,0.561}{
            \makecell[{{p{\linewidth}}}]{
                \texttt{\tiny{[P1$\rangle$GM]}}
                \texttt{förklaring: e och s är uteslutna. i, k, t finns i ordet men inte på positionerna 3, 4, 5. i kan inte vara på plats 3 (grön från gissning 1) eller plats 2 (gul här), så i måste vara på plats 1. k och t kan inte vara på plats 4 och 5 respektive. jag gissar på "ikväll" som placerar i först, k inte på 4, t inte på 5, och testar nya vanliga bokstäver v, ä, l.} \\
\texttt{gissning: ikväll} \\
            }
        }
    }
    & & \\ \\

    \theutterance \stepcounter{utterance}  
    & & & \multicolumn{2}{p{0.3\linewidth}}{
        \cellcolor[rgb]{0.9,0.9,0.9}{
            \makecell[{{p{\linewidth}}}]{
                \texttt{\tiny{[GM$|$GM]}}
                \texttt{Error: The length of the guessed word is not 5.} \\
            }
        }
    }
    & & \\ \\

    \theutterance \stepcounter{utterance}  
    & & & \multicolumn{2}{p{0.3\linewidth}}{
        \cellcolor[rgb]{0.9,0.9,0.9}{
            \makecell[{{p{\linewidth}}}]{
                \texttt{\tiny{[GM$|$GM]}}
                \texttt{game\_result = ABORT} \\
            }
        }
    }
    & & \\ \\

\end{supertabular}
}

\end{document}
